\subsection{Sprint 3}
A Sprint 3 ampliou os recursos sociais do ReadRace e trouxe o desafio de aproximar a aplicação de uma experiência utilizável fora do ambiente de desenvolvimento. O planejamento contemplava criação de clubes, chat, publicação no fórum, saída do clube e ranking completo. Também incluía configuração e resposta a quizzes, pontuação por acertos e leitura, acompanhamento de comunidades e leitores, adição de livros pelo catálogo, autenticação, personalização e disponibilização de um aplicativo instalável para Android.

No resultado coletivo, foram concluídas 16 das 17 tarefas do planejamento apresentado, além das duas tarefas de reserva: comentários em comunidades e seguir leitores. A equipe entregou fluxos de clubes, quiz e pontuação, biblioteca, login, cadastro, acesso com Google, personalização e revanche entre amigos, além de logs e alarmes de produção. Foram antecipados o acompanhamento dos desafios semanais e mensais, as listas de seguidores e seguidos e um ajuste do ranking para demonstração, inicialmente previstos para a Sprint 4. Também foram adicionados avatar e plano de fundo ao perfil. O pagamento de XP dos desafios semanais e mensais, entretanto, permaneceu para a próxima sprint.

Minha contribuição individual esteve relacionada principalmente às funcionalidades de acompanhamento de leitores, à visualização de seguidores e seguidos e ao apoio à integração das entregas do squad. Os registros pessoais de atividades fora de aula ajudam a identificar essa participação. Em 9 de outubro de 2026, registrei trabalho de desenvolvimento na issue de seguir outros usuários, associado a atividades de desenvolvimento e testes. Essa funcionalidade ampliava a relação entre os leitores para além da consulta a um perfil, oferecendo uma base para acompanhar outras pessoas dentro da plataforma.

Em 19 de outubro, trabalhei na comunicação com os colegas do squad, na resolução de conflitos de código e na revisão da issue de adição de um novo livro à biblioteca. Também realizei revisão e resolução de problemas na funcionalidade de abertura do perfil de uma comunidade e nas ações de seguir ou deixar de seguir. Essas atividades exigiram acompanhar entregas de colegas e observar como elas se conectavam aos fluxos que já existiam no aplicativo. Minha atuação, portanto, combinou implementação com revisão e apoio à organização do trabalho.

Ainda em 19 de outubro, em outro período registrado, dediquei-me ao desenvolvimento da issue de visualização de seguidos e seguidores. Essa entrega complementava a possibilidade de seguir leitores, tornando essas relações consultáveis na interface. O trabalho reforçou uma percepção que já vinha aparecendo nas sprints anteriores: funcionalidades relacionadas precisam ser tratadas como partes de uma experiência contínua, pois concluir uma ação sem oferecer uma forma clara de consultar seu resultado limita a utilidade para o usuário.

Um dos desafios coletivos da sprint foi integrar funcionalidades que alteravam as mesmas telas. A necessidade de resolver conflitos e revisar entregas relacionadas mostrou a importância de alinhar previamente os pontos de contato entre tarefas. Para mim, o aprendizado esteve na combinação entre atenção ao código e comunicação com o squad: a revisão precisava considerar também o fluxo de uso e os impactos das mudanças sobre o trabalho dos demais colegas.

Na validação do produto, foi identificado um problema nas capas dos livros no aplicativo instalado, corrigido com o uso de HTTPS. Esse episódio mostrou que o funcionamento no ambiente de desenvolvimento não garantia o mesmo comportamento em um Android real. A disponibilização completa do instalável não foi encerrada na sprint: ainda era necessário finalizar a distribuição e validar a autenticação em produção. A documentação do banco de dados também permaneceu pendente.

Essas pendências reforçaram a necessidade de distinguir implementação, publicação em produção e disponibilização ao usuário nos critérios de conclusão. Para a Sprint 4, ficaram previstos o fechamento da distribuição Android com login, o pagamento de XP dos desafios semanais e mensais, a apresentação imediata das recompensas de leitura e a aba Social com atividades dos leitores seguidos. Também permaneceram ajustes nos fluxos de clubes, biblioteca e publicações, a finalização da documentação e os testes no aplicativo instalado. Naquele planejamento, a entrega final estava prevista para 18 de novembro de 2026; os resultados dessa etapa ainda não são apresentados neste relato.
